\documentclass[conference]{IEEEtran}
\IEEEoverridecommandlockouts
% The preceding line is only needed to identify funding in the first footnote. If that is unneeded, please comment it out.
%Template version as of 6/27/2024

\usepackage{cite}
\usepackage{amsmath,amssymb,amsfonts}
\usepackage{algorithmic}
\usepackage{graphicx}
\usepackage{textcomp}
\usepackage{xcolor}
\usepackage{url}
\def\BibTeX{{\rm B\kern-.05em{\sc i\kern-.025em b}\kern-.08em
    T\kern-.1667em\lower.7ex\hbox{E}\kern-.125emX}}
\begin{document}

\title{Comparative Analysis of Algorithms for Unique Username Generation and Verification at Scale}

\author{\IEEEauthorblockN{1\textsuperscript{st} Rajesh Bhandari}
\IEEEauthorblockA{\textit{Department of Computer Engineering} \\
\textit{National College of Engineering}\\
Lalitpur, Nepal \\
nce080bct027@nce.edu.np}
\and
\IEEEauthorblockN{2\textsuperscript{nd} Utsav Bhattarai}
\IEEEauthorblockA{\textit{Department of Computer Engineering} \\
\textit{National College of Engineering}\\
Lalitpur, Nepal \\
nce080bct048@nce.edu.np}
}

\maketitle

\begin{abstract}
The widespread adoption of online platforms has created an increasing demand for unique username generation and verification systems that scale efficiently. This paper presents a comparative analysis of algorithms used for generating readable, unique usernames and verifying their availability at scale. Generation methods including trie-based suggestion, dictionary-based generation, and random suffix generation are evaluated against verification methods including Bloom filters, hashing, and exact database lookup. Algorithmic trade-offs in time complexity, space complexity, memory efficiency, accuracy, scalability, and collision handling are examined. The results identify the most balanced approach for unique username management in large-scale systems.
\end{abstract}

\begin{IEEEkeywords}
username generation, Bloom filter, trie, scalability, membership testing, algorithms
\end{IEEEkeywords}

\section{Introduction}
Today, we live in the 21st century, which is the age of technology development and digitization around the world. There are more than 8 billion people living all over the world and each year millions of new users start to use the Internet, and that's why today digital identity is extremely important.

With the growing popularity of online platforms, social networks, gaming services, and web applications, there has been a growing need for unique usernames. Usernames are used as public identifiers which make it possible to differentiate between users of the system. Moreover, usernames are essential for such processes as authentication, customization, social interaction and account management. At the same time, with the growing number of users of platforms, there arises a problem with generation of unique usernames due to the problem of name collision, username squatting and the lack of availability of common names.

To address these issues, various algorithms and data structures have been developed for generation and verification of usernames. Even though database indexing guarantees efficiency and accuracy, using the database alone can make the system inefficient due to an increased workload. That's why additional techniques such as Bloom Filters and distributed hashing mechanisms are used to improve scalability and performance in large-scale systems. Bloom filters are probabilistic algorithms which allow for quick membership tests and reduce database lookups by a huge margin \cite{b1,b2}. Similarly, consistent hashing allows for efficient mapping of user records to multiple servers and greatly reduces data shuffling when new servers join or leave the network \cite{b3}. These are some of the mechanisms that make it possible to deal with username verification more efficiently.

In case of username generation, there are a number of methods used to generate unique and user-friendly usernames. These include Random Suffix Generation where random numbers or alphanumeric characters are appended to the usernames in order to minimize collisions. There is also the Snowflake-based approach which helps in generating unique identifiers in a distributed computing environment and can be applied in the generation of usernames to guarantee uniqueness. Lastly, Trie-based approaches help in prefix search and generation of usernames based on user-specified prefixes such that the user is presented with usernames similar to what he/she wants.

The research focuses on conducting a comparative study of algorithms for generating unique usernames and verifying their uniqueness. Various methods will be examined in terms of their effectiveness, scalability, ability to ensure uniqueness, efficiency, and applicability in large-scale systems. The results of this research can be useful for choosing an appropriate solution to username management for contemporary web applications.

\section{Literature Review}
The problem of unique username generation and verification has been studied across multiple research directions, including string processing, probabilistic data structures, scalable lookup systems, and information retrieval. Although usernames appear conceptually simple from a user perspective, their management at scale introduces several algorithmic challenges related to uniqueness enforcement, low-latency availability checking, readable candidate generation, and graceful handling of name collisions when preferred identifiers are already taken \cite{b4}. Existing research consistently treats these challenges as two related but technically distinct sub-problems: generating a suitable username candidate, and verifying whether that candidate is already in use. The performance and correctness of a username system therefore depend on how effectively these two stages are designed and combined.

Research on username generation emphasizes that usability and memorability are as important as uniqueness, particularly in consumer-facing platforms where handles double as public identifiers. Studies on the structural properties of online usernames show that a large fraction of user-chosen handles are simply base names appended with numeric suffixes, a pattern that produces names which are technically unique but often unpleasant and difficult to recall \cite{b4}. The same line of work demonstrates that it is possible to generate more human-like usernames using heuristics drawn from structural onomastics, the formal study of naming patterns. By modeling usernames as structured strings rather than arbitrary character sequences, such approaches produce candidates that preserve linguistic plausibility and improve user acceptance. This motivates the study of structured name generation methods, in which usernames are constructed from patterned components rather than purely random strings, and provides a theoretical basis for evaluating generation algorithms on readability criteria rather than only on collision avoidance \cite{b4}.

Within the class of structured methods, trie-based suggestion has emerged as a particularly effective technique. A trie stores strings character by character and supports efficient prefix matching, making it well suited to autocomplete, suggestion, and incremental candidate generation tasks. When a preferred username is already taken, a trie-based system can quickly enumerate near-neighbor candidates that share a common prefix, thereby producing alternatives that remain understandable and memorable to the original user. Compared with purely random generation, trie-based suggestion yields outputs with substantially higher readability. Compared with manually specified rule-based naming, trie traversal offers better scalability because prefix operations can be performed in time proportional to the length of the input string, independent of the total size of the username corpus. However, the literature also shows that a trie alone does not guarantee uniqueness; it is a candidate-generation structure, not a verification mechanism, and must be paired with an availability check before a username can be assigned to a user \cite{b4}.

Linguistic and dictionary-based generation techniques provide an alternative to trie-driven suggestion. These methods construct usernames from curated word lists, phonetic fragments, or morphological templates, producing candidates that are pronounceable and semantically meaningful. Such approaches are especially common in practice, where many open-source libraries generate usernames by combining adjectives, nouns, and short numeric tokens drawn from a small lexicon. The principal advantage of dictionary-based generation is naturalness: outputs are easier to read, type, and remember than the unstructured strings produced by random character selection. The corresponding limitation is the bounded size of the underlying dictionary, which constrains the number of distinct candidates that can be produced before repeats become likely, and the need to recombine or extend the word list to scale to very large user populations. As a result, dictionary-based methods are frequently used in combination with random suffix generation, which appends digits or short random fragments to a base token to break ties when the dictionary alone does not suffice.

Random suffix generation is the most widely deployed fallback strategy for username creation. This technique appends digits, symbols, or short random strings to an unavailable base username, producing a candidate that is almost certainly unique with high probability. Its main advantage is simplicity: it is easy to implement, fast to execute, and imposes very little computational overhead. The corresponding disadvantage is that random suffixes degrade readability, making usernames harder to remember and less natural in appearance \cite{b4}. Within a comparative study, random suffix generation therefore represents an important baseline, as it captures the trade-off between low implementation complexity and reduced human usability, and serves as a useful lower bound against which more sophisticated methods can be measured.

Complementing generation, the verification stage of a username system has been extensively studied in the context of distributed authentication and large-scale membership testing. Bloom filters, which are space-efficient probabilistic data structures that support set membership queries, have received particular attention. Prior work on user identifier lookup in distributed authentication systems shows that Bloom filters can substantially reduce authentication latency and network traffic by performing rapid preliminary checks for identifier existence before consulting the main database \cite{b5}. The same principle applies directly to username verification: a Bloom filter can quickly determine that a username is definitely not present, thereby eliminating many unnecessary database round-trips. This property is especially valuable in large-scale systems, where availability checks occur at high frequency and database access dominates end-to-end latency.

The literature also emphasizes that Bloom filters are probabilistic rather than exact. A Bloom filter can produce a false positive, indicating that a username might already be taken even when it is in fact available, but under standard operating conditions it does not produce false negatives. Because of this asymmetry, Bloom filters are best deployed as a pre-filter in front of an authoritative store, not as the final source of truth. The exact decision must still be enforced by a deterministic database lookup, a unique constraint in persistent storage, or a comparable exact membership mechanism. This layered design pattern, in which a fast probabilistic check filters out impossible candidates and an exact check resolves the remainder, has been shown to substantially improve throughput while preserving correctness, and it appears consistently across the literature on scalable membership testing \cite{b5}.

Research on scalable name lookup in networking and content-centric architectures further reinforces this design pattern. A two-stage Bloom filter scheme applied to name lookup in Named Data Networking has been shown to significantly improve lookup throughput while using less memory than traditional character-trie structures, demonstrating that layered probabilistic filtering can outperform single-structure approaches at scale \cite{b6}. Although this work is not directly about usernames, the underlying problem, namely efficient search over very large sets of strings with bounded memory, is structurally similar. The study demonstrates that Bloom filters can be combined with structural lookup strategies, such as tries and prefix-indexed data structures, to achieve high performance without sacrificing memory efficiency, and provides a strong conceptual basis for hybrid username verification systems that combine a fast pre-check with an exact database validation.

Hybrid deterministic--probabilistic lookup structures constitute a third line of relevant research. These systems integrate prefix trees with probabilistic filters so that a Bloom filter pre-search stage rapidly eliminates the majority of trie nodes that cannot contain a match, after which a small number of full trie traversals complete the lookup. The result is a search architecture that benefits from the prefix-pruning efficiency of the trie and the constant-time membership testing of the Bloom filter, while keeping memory usage manageable through sharing and partitioning. In the context of username systems, this hybrid pattern is particularly relevant because trie-based candidate generation and Bloom-filter-assisted verification can be viewed as complementary operations on the same data structure: the trie helps generate readable candidate names or prefixes, while the Bloom filter reduces unnecessary database access during verification \cite{b6,b7,b8}.

Beyond Bloom filters and tries, alternative verification approaches based on hashing and direct database lookup are also widely used. Hashing methods map usernames to fixed-size digests stored in a hash table, providing expected constant-time lookup at the cost of collision handling. Exact database lookup, typically backed by a B-tree or hash index on a unique column, provides guaranteed correctness but incurs higher per-query latency, particularly in distributed deployments. Multi-stage filtering architectures that combine Bloom filters, caches, and indexed databases have been proposed to balance these trade-offs, eliminating obviously unavailable candidates with a Bloom filter, serving popular lookups from cache, and reserving exact database access for the residual ambiguous cases. Adaptive prefix Bloom filter designs further refine this idea by dynamically adjusting the depth of probabilistic prefix matching to current workload conditions, yielding a lookup scheme that scales gracefully with corpus size while maintaining low false-positive rates \cite{b9,b10}.

A consistent observation across these studies is that no single algorithm simultaneously satisfies the full set of requirements for unique username management at scale. Structured and trie-based generation methods produce readable candidates but cannot, by themselves, guarantee uniqueness. Random suffix generation provides a simple uniqueness mechanism but degrades usability. Bloom filters offer fast probabilistic verification but cannot replace exact checks. Hashing and indexed database lookup guarantee correctness but become latency bottlenecks under heavy load. Each technique therefore embodies a distinct trade-off among time complexity, space complexity, memory efficiency, accuracy, scalability, collision handling, and practical usability. The literature has repeatedly converged on the conclusion that performance and correctness are achieved not by any single algorithm, but by combining complementary methods into a layered pipeline in which the strengths of one stage compensate for the limitations of another.

A further limitation of existing work is that the candidate-generation and verification stages are typically studied in isolation. Research on username structure focuses on readability and linguistic plausibility but rarely evaluates verification cost at scale, while research on Bloom filters and scalable lookup focuses on membership-test efficiency but typically assumes a fixed or randomly distributed set of keys that does not reflect the heavy-tailed, prefix-correlated distribution of real usernames. As a result, no prior study provides a unified empirical comparison of generation and verification algorithms under realistic large-scale conditions, and it remains unclear how the algorithmic trade-offs observed in each individual study manifest when the methods are combined in a complete username management pipeline. This gap motivates a comparative analysis that evaluates trie-based suggestion, dictionary-based generation, and random suffix generation against Bloom-filter-assisted verification, hash-based lookup, and exact database access within a single experimental framework, measuring not only speed and memory but also collision rates, false-positive rates, and the readability of generated candidates. The methodology and experimental design that follow from these observations are presented in the next section.

\section*{References}
\begin{thebibliography}{00}
\bibitem{b1} L. Luo, D. Guo, R. T. B. Ma, O. Rottenstreich, and X. Luo, ``Optimizing Bloom Filter: Challenges, Solutions, and Comparisons,'' \emph{arXiv preprint arXiv:1804.04777}, 2018.
\bibitem{b2} L. Ramabaja and A. Avdullahu, ``The Distributed Bloom Filter,'' \emph{arXiv preprint arXiv:1910.07782}, 2019.
\bibitem{b3} D. Karger, E. Lehman, T. Leighton, M. Levine, D. Lewin, and R. Panigrahy, ``Consistent Hashing and Random Trees: Distributed Caching Protocols for Relieving Hot Spots on the World Wide Web,'' in \emph{Proc. 29th Annual ACM Symposium on Theory of Computing (STOC)}, El Paso, TX, USA, 1997, pp. 654--663.
\bibitem{b4} ``Structural Onomatology for Username Generation: A Partial Account,'' Universitat de Lisboa / University of Minho, 2020. [Online]. Available: \url{https://repositorium.sdum.uminho.pt/bitstream/1822/69231/1/STAIRS20_P26_paper.pdf}
\bibitem{b5} ``Fast and Efficient UserID Lookup in Distributed Authentication: A Probabilistic Approach Using Bloom Filters,'' \emph{International Journal of Computer Engineering}, vol. 6, no. 2, pp. 1--16, 2024. [Online]. Available: \url{https://carijournals.org/journals/IJCE/article/download/2124/2512/6389}
\bibitem{b6} Y. Wang, T. Pan, Z. Mi, H. Dai, X. Guo, T. Zhang, B. Liu, and Q. Dong, ``NameFilter: Achieving fast name lookup with low memory cost via applying two-stage Bloom filters,'' in \emph{Proceedings of IEEE INFOCOM}, 2013, pp. 95--99.
\bibitem{b7} G. Holley, A. Antonini, et al., ``Dual-load Bloom filter: Application for name lookup,'' \emph{Computer Communications}, 2019.
\bibitem{b8} ``Name Prefix Matching Using Bloom Filter Pre-Searching,'' in \emph{Proceedings of the Eleventh ACM/IEEE Symposium on Architectures for Networking and Communications Systems}, 2015.
\bibitem{b9} ``A New Bloom Filter Architecture for FIB Lookup in Named Data Networking,'' 2024.
\bibitem{b10} ``Scalable Name Lookup with Adaptive Prefix Bloom Filter for Named Data Networking,'' 2013.
\end{thebibliography}

\end{document}
